% =====================================================================
\part{The equilibrium surface: log-odds coordinates, branches, and solvers}
% =====================================================================

\chapter{The log-odds configuration}

\section{Specification}
The final coordinate system, run at the user's direction:
\begin{itemize}\itemsep2pt
  \item \textbf{$u$-grid}: uniform in the no-learning private posterior
        log-odds transform $\xi = \tanh(\tau u / 2)$,
        $\xi \in [-0.85, 0.85]$, mapped back by
        $u = (2/\tau)\operatorname{atanh}(\xi)$. No bounded cube: the
        grid self-scales to $\approx \pm 5\sigma_{\rm signal}$ at every
        $\tau$ (at $\tau = 0.005$ the cube spans $\pm 70$ in $u$; at
        $\tau = 1$, $\pm 2.4$).
  \item \textbf{$p$-grid}: uniform in public-only posterior log-odds
        $\Lpub(p)$, bootstrap-computed from the cube empirical CDF and
        \emph{antisymmetrized by construction}
        ($\Lpub \mapsto \tfrac12(\Lpub - \Lpub^{\rm rev})$) after the
        unsymmetrized bootstrap was caught breaking the equilibrium's
        sign-flip symmetry at $O(0.4)$.
  \item \textbf{Operator}: \strict{} co-area + POU (zero kernel).
  \item \textbf{Clearing}: exact CRRA bisection (the linearized
        clearing bug and its repair are recorded in Part I).
  \item \textbf{Residual}: both raw $|F|_\infty$ and the
        density-weighted $\Fw = \max |\Phi(P) - P| \cdot W_3$,
        $W_3 \propto \prod_i \bar f(u_i)$ --- the raw max over an
        unbounded cube is dominated by corner cells of vanishing
        probability and is not the economically meaningful residual.
  \item \textbf{Symmetry projection} (permutation + sign-flip) every
        iteration; live per-iteration reporting.
\end{itemize}

\section{The $G$-ladder and $\gamma$-sweep}
At $(\gamma, \tau) = (100, 0.1)$: $\Fw$ best of $1.3\times10^{-2}$
($G{=}7$) down to $1.6\times10^{-3}$ ($G{=}9$), non-monotone across
$G \in \{7..15\}$ --- the \strict{} stall manifests differently per
grid. The $\gamma$-sweep at $G = 11$ spans
$\Fw = 1.8\times10^{-2}$--$2.6\times10^{-2}$ across
$\gamma \in [1, 1000]$: genuinely $\gamma$-dependent (after the clearing
fix), all cells with sign-flip symmetry at $1.1\times10^{-16}$ and
$P(\mathbf{0}) = \tfrac12$ exactly.

\begin{figure}[H]\centering
  \includegraphics[width=0.9\linewidth]{lo_figs/g100_t0.1000_G11.png}
  \caption{The log-odds configuration at
    $(\gamma, \tau, G) = (100, 0.1, 11)$: raw vs density-weighted
    residual traces (left) and the equilibrium slice (right). The raw
    residual is dominated by zero-density corners; the weighted
    residual is the meaningful one.}
\end{figure}

\section{The 50-cell $\tau$-ladder}
Fifty log-spaced $\tau \in [0.001, 1.2]$, chained warm-start,
16 seconds total. Three regimes:
\[
  \underbrace{\tau \le 0.005}_{\Fw \approx 2.5\times10^{-15}}
  \quad\longrightarrow\quad
  \underbrace{0.005 \lesssim \tau \lesssim 0.02}_{\Fw: 10^{-7} \to 10^{-3}}
  \quad\longrightarrow\quad
  \underbrace{\tau \gtrsim 0.05}_{\Fw \approx 10^{-2} \text{ (stall)}}
\]
with perfect sign-flip symmetry everywhere and the price range opening
monotonically with $\tau$.

\begin{figure}[H]\centering
  \includegraphics[width=\linewidth]{lol_figs/ladder_summary.png}
  \caption{The 50-cell Picard $\tau$-ladder: residual regimes (left) and
    equilibrium price range vs $\tau$ (right).}
\end{figure}

\begin{figure}[H]\centering
  \includegraphics[width=\linewidth]{lol_figs/logodds_contours.png}
  \caption{Equilibrium log-odds contours
    $\operatorname{logit}P(u_1, u_2 \mid u_3{=}0)$ on private-log-odds
    axes across six $\tau$ values: a clean diagonal zero contour at low
    $\tau$ deforming into the kinked high-$\tau$ structure whose
    critical values generate the lookup cusps.}
\end{figure}

\begin{figure}[H]\centering
  \includegraphics[width=\linewidth]{lol_figs/logodds_contours_u.png}
  \caption{The same surfaces on raw $u$ axes with $k\sigma$ density
    rings: the unbounded-cube grid keeps the Gaussian mass interior at
    every $\tau$.}
\end{figure}

\section{Anatomy of one lookup evaluation}
For each $(p, u_k)$ query the \strict{} builder consumes: all $G^2$
slice vertex values (linear interpolation backbone), $2 \times 16$
Gauss--Legendre sweep nodes, and the level-set roots found per sweep
line --- typically crossing 25--56 of the $(G{-}1)^2$ cells. One full
table is $G_p \times G = 891$ such queries.

\begin{figure}[H]\centering
  \includegraphics[width=\linewidth]{lol_figs/lookup_points.png}
  \caption{Which grid points the lookup uses: cube vertices (black), GL
    sweep lines (light), level-set roots from each sweep (blue
    $\times$ / orange $+$), and the cells crossed by the level set
    (red). Top row: $\tau = 0.005$ --- the level set hugs the
    antidiagonal and off-median prices find no roots (the fallback
    region of Finding 14). Bottom: $\tau = 0.5$ --- well-resolved level
    sets across the cube.}
\end{figure}

\chapter{Branches and Newton-type solvers}

\section{The trivial branch is exact}
$P \equiv \tfrac12$ satisfies the K=3 CRRA clearing identically (all
posteriors equal the prior, all $R_i = 1$, zero demand): cold-started
solvers of every type converge to it instantly, to $F = 0$ exactly.
The equilibrium correspondence is multi-valued everywhere; \emph{branch
identification is part of the computation's burden of proof}. The
informative branch is reached only by continuation (warm-starting along
the $\tau$-ladder from the analytic low-$\tau$ solution).

\section{Newton--Krylov via the implicit function theorem}
The IFT Newton direction
$\delta P = -(\partial\Phi/\partial P - I)^{-1}(\Phi(P) - P)$,
implemented matrix-free by LGMRES with finite-difference Jacobian-vector
products and Wolfe line search, warm-started from the Picard ladder:
\begin{itemize}\itemsep2pt
  \item \textbf{Low $\tau$}: weighted residuals driven as deep as
        $10^{-25}$--$10^{-50}$ on cells where the basin is locally
        smooth --- far below the Picard floor.
  \item \textbf{High $\tau$}: plateaus at $\Fw \sim 10^{-3}$--$10^{-2}$
        with intermittent symmetry breaks --- the finite-difference
        Jacobian inherits the operator's $C^0$-not-$C^1$ roughness.
  \item \textbf{Cost}: 200--1400 $F$-evaluations per cell
        ($50$--$100\times$ Picard); total ladder 207\,s vs 16\,s.
\end{itemize}
MINPACK HYBRD (explicit numerical Jacobian + Broyden updates) was also
tried: in logit space it is captured by the trivial branch even from
informative warm starts; in $P$-space it exits the unit cube
(one cell reached $P \in [-1590, 1591]$) even with bound penalties.
Newton--Krylov with continuation warm starts is the only stable Newton
variant found.

\begin{figure}[H]\centering
  \includegraphics[width=\linewidth]{lon_figs/ladder_newton.png}
  \caption{The IFT-Newton ladder vs Picard: deeper convergence at low
    $\tau$, the same operator wall at high $\tau$, at $50$--$100\times$
    the cost. Newton refines; it does not rescue.}
\end{figure}

\section{Synthesis of the surface-side evidence}
Across Picard, Anderson, Newton--Krylov, HYBRD, smoothed-$\varepsilon$
operators, frozen-knot inner loops, and the structural ansatz: every
method that respects the \strict{} operator's discretization plateaus at
the same $G$-dependent wall at high $\tau$, and every method that
smooths past the wall pays in bias. Combined with the
$C^0$-not-$C^1$ theorem and the success of the consistent-kernel joint
limit at $10^{-14}$ (\texttt{k3\_coarea\_limit}), the conclusion is
structural: \textbf{the high-resolution surface computation must use a
consistent $C^\infty$ operator (kernel co-area in the central-path
limit), with the bare \strict{} operator reserved for the role of
unbiased cross-check at matched tolerance}. That division is the spine
of the recommendation in Part VII.
